\documentclass{article}
\usepackage{amsmath, amssymb}
\usepackage{geometry}
\usepackage{booktabs}
\geometry{margin=2.5cm}

\title{Game Analysis: The (10, 10) Rock Game}
\author{}
\date{}

\begin{document}
\maketitle

\noindent Two piles of 10 rocks each. On your turn, take 1 from Pile A, 1 from Pile B, or 1 from both. The player who takes the last rock wins. Your opponent moves first.

\section*{W-positions and L-positions}

A position is a \textbf{L-position} if the player about to move loses (with optimal play), and an \textbf{W-position} if they win. The terminal position $(0,0)$ is a P-position since the player facing it has already lost.

\section*{The Pattern}

Seeing the pattern in small cases reveals:

\[
(a, b) \text{ is a L-position} \iff a \equiv 0 \pmod{2} \text{ and } b \equiv 0 \pmod{2}
\]

\textbf{Proof:}

\medskip
\noindent\textit{Both even $\Rightarrow$ every move leads to an W-position.}
The three moves give $(a-1,b)$, $(a,b-1)$, $(a-1,b-1)$. In each case at least one coordinate becomes odd, so no move reaches a L-position. Hence the current player loses. 

\medskip
\noindent\textit{At least one odd $\Rightarrow$ some move leads to a L-position.}
\begin{itemize}
  \item $a$ odd, $b$ even: take from A $\to$ $(a-1, b)$, both even.
  \item $a$ even, $b$ odd: take from B $\to$ $(a, b-1)$, both even.
  \item Both odd: take from both $\to$ $(a-1, b-1)$, both even.
\end{itemize}
In every case there is a move to a L-position. Hence the current player wins. 

\section*{Result for (10, 10)}

Since $10 \equiv 0 \pmod{2}$, the position $(10,10)$ is a \textbf{L-position}.The opponent moves first and faces a L-position, so \textbf{we win} with optimal play.

\section*{Winning Strategy}

After every opponent move, restore both piles to even numbers:

\begin{itemize}
  \item Takes from A only ($a$ odd, $b$ even): take 1 from A $\to$ both even.
  \item Takes from B only ($a$ even, $b$ odd): take 1 from B $\to$ both even.
  \item Takes from both (both odd): take from both $\to$ both even.
\end{itemize}

\noindent We always hand our opponent a L-position. This continues until you take the last rock(s) at $(2,2) \to (1,1) \to (0,0)$ and win.

\end{document}